\section*{Capítulo \ref{ch:g12:mountain-building} --- A formação das montanhas}

\begin{solution}{exo:g12:mountain-building:1}
(a) \omterm{def:g12:mountain-building:thrust}{cavalgamento}; (b) \omterm{def:g12:mountain-building:ophiolite}{ofiolito}; (c) \omterm{def:g12:mountain-building:fold}{sinclinal}; (d) \omterm{def:g12:mountain-building:suture}{sutura}; (e) \omterm{def:g12:mountain-building:fold}{anticlinal};
(f) \omterm{def:g12:mountain-building:nappe}{nappe}.
\end{solution}

\begin{solution}{exo:g12:mountain-building:2}
(1) Os flancos mergulham para longe do eixo da \omterm{def:g12:mountain-building:fold}{dobra}: um \omterm{def:g12:mountain-building:fold}{anticlinal}. (2) As
camadas mais antigas. (3) Não: a erosão segue as camadas duras e moles, não
a forma da \omterm{def:g12:mountain-building:fold}{dobra}; o núcleo de um \omterm{def:g12:mountain-building:fold}{anticlinal} pode ser escavado num vale.
\end{solution}

\begin{solution}{exo:g12:mountain-building:3}
(1) $e = 1 - \cos\ang{30} = 1 - 0.866 = 0.134$, cerca de \qty{13}{\percent}.
(2) $L_0 = 6.0/\cos\ang{30} \approx \qty{6.9}{km}$.
\end{solution}

\begin{solution}{exo:g12:mountain-building:4}
(1) Uma \omterm{def:g11:ocean-ridges:normal-fault}{falha normal}; a crosta foi estirada. (2) Sim: é uma \omterm{def:g12:mountain-building:thrust}{falha inversa}
que mergulha menos de \ang{45}.
\end{solution}

\begin{solution}{exo:g12:mountain-building:5}
(1) Sedimentos marinhos, basaltos em almofada, complexo de \omterm{def:g11:igneous-rocks:intrusion}{diques}, \omterm{def:g11:ocean-ridges:gabbro}{gabro},
\omterm{def:g11:ocean-ridges:peridotite}{peridotito}. (2) A lava só forma almofadas quando esfria na água. (3) Um
\omterm{def:g12:mountain-building:ophiolite}{ofiolito} é um pedaço do oceano que ficava entre dois continentes; quando eles
colidiram, lascas do seu fundo foram empurradas sobre as suas bordas em vez
de afundar, por isso elas se encontram onde o oceano se fechou, ao longo da
\omterm{def:g12:mountain-building:suture}{sutura}.
\end{solution}

\begin{solution}{exo:g12:mountain-building:6}
(1) $L_0 = 2.1 + 1.8 + 2.4 + 1.6 = \qty{7.9}{km}$;
$e = (7.9 - 6.0)/7.9 \approx 0.24$, \qty{24}{\percent}. (2) $\cos\theta =
6.0/7.9 \approx 0.759$, $\theta \approx \ang{41}$.
\end{solution}

\begin{solution}{exo:g12:mountain-building:7}
(1) $45\cos\ang{10} \approx \qty{44.3}{km}$. (2) $45\sin\ang{10} \approx
\qty{7.8}{km}$. (3) \omterm{def:g12:mountain-building:fold}{Dobras} assim encurtam de $1 - \cos\ang{12} \approx 0.0219$:
o cinturão precisaria ter $44.3/0.0219 \approx \qty{2000}{km}$ de largura
antes do dobramento.
\end{solution}

\begin{solution}{exo:g12:mountain-building:8}
(1) A superposição: lida sem a falha, a rocha de cima seria tomada pela mais
jovem. (2) $500 - 75 = 425$ milhões de anos. (3) $\qty{40000}{m} /
\qty{800}{m} = 50$ vezes.
\end{solution}

\begin{solution}{exo:g12:mountain-building:9}
(1) $(30.23 - 19.38) \times 111.2 = 10.85 \times 111.2 \approx
\qty{1207}{km}$. (2) $1207/14 \approx \qty{86}{km}$ por milhão de anos,
\qty{86}{mm} ou cerca de \qty{8.6}{cm} por ano. (3) \qty{9}{m} por século
são \qty{9}{cm} por ano: os dois concordam dentro de cerca de
\qty{5}{\percent}.
\end{solution}

\begin{solution}{exo:g12:mountain-building:10}
(1) $H = 35/(1 - 0.40) \approx \qty{58}{km}$. (2) $\Delta H \approx
\qty{23.3}{km}$; ela fica mais alta de $23.3 \times 550/3380 \approx
\qty{3.8}{km}$, e a sua \omterm{def:g12:continental-crust:root}{raiz} a mais tem $23.3 - 3.8 \approx \qty{19.5}{km}$ de
profundidade.
\end{solution}

\begin{solution}{exo:g12:mountain-building:11}
(1) O desenho multiplica $\tan\theta$ por 3: $\tan\theta = \tan\ang{40}/3
\approx 0.280$, $\theta \approx \ang{15.6}$. (2) Verdadeiro: $1 - \cos\ang{15.6}
\approx 0.037$, \qty{3.7}{\percent}; lido no desenho: $1 - \cos\ang{40}
\approx 0.234$, \qty{23}{\percent}. (3) Uma escala vertical esticada torna
os mergulhos mais íngremes e alonga as camadas, e assim faz o \omterm{def:g12:mountain-building:shortening}{encurtamento}
parecer várias vezes grande demais.
\end{solution}

\begin{solution}{exo:g12:mountain-building:12}
(1) $L_0 = 10/(1 - 0.317) \approx \qty{14.6}{km}$. (2) $1 - \cos\ang{57.5}
\approx 1 - 0.537 = 0.463$, \qty{46}{\percent}. (3) Na \omterm{def:g12:mountain-building:fold}{dobra} arredondada,
só uma pequena parte da camada mergulha a \ang{57.5}; perto das charneiras
ela é quase horizontal, e as partes horizontais não são encurtadas. Na \omterm{def:g12:mountain-building:fold}{dobra}
de flancos retos, a camada inteira mergulha a \ang{57.5}.
\end{solution}

\begin{solution}{exo:g12:mountain-building:13}
(1) $\sqrt{16.6^2 + 41.8^2} \approx \qty{45.0}{mm/yr}$. (2) $41.8 - 23.5 =
\qty{18.3}{mm}$ por ano, isto é, cerca de \qty{18}{km} de \omterm{def:g12:mountain-building:shortening}{encurtamento}
norte--sul por milhão de anos ao norte de Lhasa. (3) A própria Lhasa se
moveu para o norte, e os terremotos do cinturão se espalham muito ao norte
do Himalaia, através da Ásia: a crosta ali ainda está se rompendo.
\end{solution}

\begin{solution}{exo:g12:mountain-building:14}
(1) O seu mergulho é próximo de zero, logo o \omterm{def:g12:mountain-building:shortening}{encurtamento} é o deslocamento,
\qty{45}{km}. (2) $45/(1 - \cos\ang{12}) = 45/0.0219 \approx \qty{2060}{km}$.
(3) As grandes convergências são absorvidas principalmente pelos
\omterm{def:g12:mountain-building:thrust}{cavalgamentos} e pelas \omterm{def:g12:mountain-building:nappe}{nappes} e pelas \omterm{def:g12:mountain-building:fold}{dobras} apertadas; as \omterm{def:g12:mountain-building:fold}{dobras} suaves
absorvem muito pouco.
\end{solution}

\begin{solution}{exo:g12:mountain-building:15}
(1) $\qty{8848.86}{m} = \qty{8848860}{mm}$ em $\num{5e7}$ anos: cerca de
\qty{0.18}{mm} por ano. (2) $40/0.18 \approx 220$: a subida é mais de
duzentas vezes mais lenta que o movimento da Índia. (3) O calcário foi
depositado abaixo do nível do mar, e a erosão removeu as rochas que ficavam
acima dele: ele foi erguido mais do que a sua altitude atual.
\end{solution}

\begin{solution}{pb:g12:mountain-building:1}
\textbf{1.} $2\pi \times 6371/360 \approx \qty{111.2}{km}$.

\textbf{2.} Entre há 50 milhões de anos (\ang{0.68}~Sul) e há 45
(\ang{3.30}~Norte): há cerca de $50 - 5 \times 0.68/3.98 \approx 49$ milhões
de anos.

\textbf{3.} $(19.38 - 0.68) \times 111.2 = 18.70 \times 111.2 \approx
\qty{2079}{km}$.

\textbf{4.} $2079/16 \approx \qty{130}{km}$ por milhão de anos: \qty{130}{mm}
por ano, cerca de \qty{13}{cm}.

\textbf{5.} 66 a 60: $7.28 \times 111.2/6 \approx 135$; 60 a 55: $4.86
\times 111.2/5 \approx 108$; 55 a 50: $6.56 \times 111.2/5 \approx 146$; 50
a 45: $3.98 \times 111.2/5 \approx 89$; 45 a 40: $2.29 \times 111.2/5
\approx 51$ (\unit{mm} por ano). Mais rápida de 55 a 50 milhões de anos atrás.

\textbf{6.} $(20.00 - 5.59) \times 111.2/40 = 14.41 \times 111.2/40 \approx
\qty{40}{mm}$ por ano.

\textbf{7.} $40/130 \approx 0.31$: menos de um terço.

\textbf{8.} Do passo 55--50 (146) para o passo 50--45 (89): por volta de 50
milhões de anos atrás.

\textbf{9.} A \omterm{def:g12:continental-crust:continental-crust}{crosta continental} é leve demais para afundar: quando ela
chegou à fossa, a \omterm{def:g11:subduction-zones:subduction}{subducção} não podia mais puxar a Índia livremente, e o
outro continente a freou.

\textbf{10.} $(20.00 + 0.68) \times 111.2 = 20.68 \times 111.2 \approx
\qty{2300}{km}$.

\textbf{11.} $(29.65 - 21.15) \times 111.2 = 8.50 \times 111.2 \approx
\qty{945}{km}$.

\textbf{12.} $2300 - 945 \approx \qty{1354}{km}$.

\textbf{13.} Hoje $(29.65 - 20.00) \times 111.2 \approx \qty{1073}{km}$;
há 50 milhões de anos $(21.15 + 0.68) \times 111.2 = 21.83 \times 111.2
\approx \qty{2427}{km}$ (e $2427 - 1073 = 1354$).

\textbf{14.} $e = 1354/2427 \approx 0.56$, \qty{56}{\percent}.

\textbf{15.} \omterm{def:g12:mountain-building:fold}{Dobras}, e \omterm{def:g12:mountain-building:thrust}{falhas inversas}: \omterm{def:g12:mountain-building:thrust}{cavalgamentos} que levam \omterm{def:g12:mountain-building:nappe}{nappes}.

\textbf{16.} $H = 35 \times 2427/1073 \approx \qty{79}{km}$.

\textbf{17.} $\Delta H = 79 - 35 = \qty{44}{km}$; ela ficaria mais alta de
$44 \times 550/3380 \approx \qty{7.2}{km}$.

\textbf{18.} $\Delta H = 4.6 \times 3380/550 \approx \qty{28}{km}$, uma
crosta de $35 + 28 = \qty{63}{km}$ de espessura, perto dos \qty{70}{km}
medidos.

\textbf{19.} A crosta não conserva a sua seção: partes da Ásia são
espremidas para o leste, e a erosão remove rocha do topo das montanhas.

\textbf{20.} $2079 + 2300 \approx \qty{4380}{km}$: desde o fim do
Cretáceo, a Índia viajou cerca de \textbf{\qty{4400}{km} para o norte}.
\end{solution}
